\begin{table}[htbp]
\centering
\footnotesize
\caption{Selected prior work and the space occupied by this study. The table is a
focused comparison, not a systematic review. It includes studies that combine
Jakarta relevance, multi-indicator heat prioritisation, accessibility or explicit
ranking-disagreement analysis. The individual ingredients of the Jakarta index
have precedent; the contribution claimed here is the attribution and uncertainty
analysis, not the combination itself.}
\label{tab:literature}
\begin{threeparttable}
\begin{tabularx}{\textwidth}{@{}>{\raggedright\arraybackslash}p{2.35cm}
>{\raggedright\arraybackslash}p{2.45cm}
>{\raggedright\arraybackslash}p{2.65cm}
>{\raggedright\arraybackslash}p{2.55cm}
>{\raggedright\arraybackslash}p{2.25cm}
>{\raggedright\arraybackslash}X@{}}
\toprule
Study & Context / unit & Indicators / construct & Spatial method / uncertainty & Closest overlap & Evidence-backed distinction \\
\midrule
Perdinan et al.\ (2026) & Jakarta; 30~m pixels, municipal aggregation & LST, vegetation deficit, built-up intensity, population density & OLS regression; weighted priority index; four weight scenarios and threshold relaxation & Jakarta multi-indicator adaptation prioritisation and stability testing & Regression-informed adaptation tiers, not our sequential screening models; no age count/share, candidate-park network uncertainty or Shapley decomposition \\
Fitria et al.\ (2026) & Taiwan; 158 districts & 17 exposure, sensitivity and adaptive-capacity indicators & Equal-weighting versus PCA; Spatial Disagreement Index; leave-one-out sensitivity & Explicit ranking disagreement and decision support & Disagreement between aggregation methods; we decompose a sequential construct transition and modelled candidate-park uncertainty \\
Ramadhanis et al.\ (2026) & Jakarta; fishnet grid & Under-five population density; OSM green-space distance & Network shortest path; child green-space vulnerability index & Jakarta children and green-space access & Age-specific vulnerability map; we compare age counts with shares in a multi-model ranking and use all-age candidate parks \\
Dong et al.\ (2024) & Zhengzhou; 2010, 2015 and 2020 & Urban green space, heat-island area, residents & Improved 2SFCA; bivariate local Moran's $I$ & Cooling accessibility combined with heat-risk spatial statistics & Accessibility is the estimand; no M1--M4 composite ranking or attribution decomposition \\
Elmarakby \& Elkadi (2024) & Manchester; districts & UHI intensity, green cover, population density & Heat Risk Index; correlation; planning tiers & Composite heat-risk prioritisation & District risk index, not cell-level screening; no age count/share or network uncertainty \\
Ulfiasari \& Yola (2022) & Jakarta metropolitan area & Landsat UHI intensity and area over ten years & Remote-sensing change mapping; planning interpretation & Jakarta UHI mapping for spatial planning & Maps UHI distribution and trends; no multi-indicator priority ranking or rank-change analysis \\
Kim \& Kwak (2026) & Older adults; walking network & Cooling-weighted park accessibility & Generalised two-step floating catchment area & Age and park-access measurement & Focus on the accessibility model rather than Jakarta ranking sensitivity or count/share redundancy \\
This study & Five mainland Jakarta cities; 500~m grid & LST, vegetation, population, age count/share, candidate-park distance & M1--M4; Jaccard; Shapley; 52 scenarios; LISA & Attribution, demographic redundancy and candidate-park uncertainty & Conditional screening analysis; not a validated risk, health or public-access model \\
\bottomrule
\end{tabularx}
\begin{tablenotes}[flushleft]\footnotesize
\item LST: land surface temperature; UHI: urban heat island; 2SFCA: two-step floating catchment area; PCA: principal component analysis. The table summarises the research question and evidence used in each study; it does not claim to reproduce every method. The absence of a study is not evidence that no similar work exists.
\end{tablenotes}
\end{threeparttable}
\end{table}
